\openleft\chapter{Summary}

The landscape of theoretical condensed matter physics has transformed dramatically over the past century, with computers now playing a central role in scientific discovery. In this thesis, we apply computer science techniques to simplify, speed up, and improve quantum many-body simulations.

The challenge in quantum many-body simulations arises from a fundamental difference between classical and quantum systems. In classical mechanics, describing the motion of a pool ball requires storing just 2 coordinates (for example, $x$ and $y$ positions, assuming the ball moves only horizontally). For 3 balls, we need $2+2+2=6$ coordinates; for 10 balls, we need 20 coordinates. The number of coordinates grows linearly with the number of particles, making it straightforward to simulate thousands of particles on a computer.

Quantum physics is different. If we need 2 numbers to describe the state of one quantum particle, then, in general, describing 3 such particles requires $2\cdot 2\cdot 2=8$ numbers, because the state of one particle ``depends'' on the states of all other particles. Describing 10 particles then requires $2^{10}=1024$ numbers. This exponential scaling is fundamental to quantum many-body theory and makes simulating even 100 particles extremely challenging, as one would need to store $1267650600228229401496703205376$ numbers.

Despite this exponential complexity, there are several limiting cases where we can hope to solve the problem. Chapters~\ref{ch:prb18} and~\ref{ch:mat22} focus on weakly interacting particles, where the required storage reduces to the same linear scaling as in classical mechanics. For strongly interacting particles, we can attempt exact solutions but only for small systems. This approach is discussed in Chapters~\ref{ch:ls21},~\ref{ch:paw-atm}, and~\ref{ch:corr-hops}. Finally, modern machine learning methods offer approximate solutions for strongly interacting cases, though they encounter an analogue of the so-called sign problem. Chapters~\ref{ch:nc20} and~\ref{ch:cp23} examine this problem in depth and introduce new algorithms to address it.

The techniques developed in this work advance our computational capabilities in quantum many-body theory and demonstrate how interdisciplinary approaches can tackle challenging numerical problems.
% The techniques developed in this work do not solve all problems in quantum many-body theory, but they provide a foundation for future method development and demonstrate the power of interdisciplinary approaches combining physics and computer science.
% This thesis presents a few concrete examples how a combination of computer science and physics techniques can help tackle challenges in quantum many-body theory. I hope that the tools developed in this work provide a foundation for future methodological developments.
